\documentclass[10pt,a4paper]{amsart}
\usepackage{amsmath,amssymb,amsthm,graphicx}
% Retained preamble for the Stacks Project chapter {CH} (shared project preamble).
% The chapter file pinned-original-source.tex only carries \input{{preamble}} on its first line,
% so no preamble text lives inside the chapter itself.  The Stacks Project uses one single shared
% preamble for every chapter: the snapshot stages it three times (sources/stacks-chow/preamble.tex,
% sources/stacks-etale-cohomology/preamble.tex and focused-seeds/AGP/sources/stacks-snapshot-20261011/
% preamble.tex) and those copies are identical in size.  The block below is copied mechanically,
% line for line, from sources/stacks-chow/preamble.tex lines 14-18 and 154-217: the Xy-pic setup and
% the native theorem declarations together with the mathematical macros.  Nothing was added, removed
% or redefined; the wrapper class and the amsmath/amssymb/amsthm/graphicx base packages are supplied
% by the assembler.
% For commutative diagrams we use Xy-pic
% Chapter: stacks-spaces-more-morphisms.tex -- More on Morphisms of Spaces.
\usepackage[all]{xy}

% We use 2cell for 2-commutative diagrams.
\xyoption{2cell}

% Theorem environments.
%
\theoremstyle{plain}
\newtheorem{theorem}[subsection]{Theorem}
\newtheorem{proposition}[subsection]{Proposition}
\newtheorem{lemma}[subsection]{Lemma}

\theoremstyle{definition}
\newtheorem{definition}[subsection]{Definition}
\newtheorem{example}[subsection]{Example}
\newtheorem{exercise}[subsection]{Exercise}
\newtheorem{situation}[subsection]{Situation}

\theoremstyle{remark}
\newtheorem{remark}[subsection]{Remark}
\newtheorem{remarks}[subsection]{Remarks}

\numberwithin{equation}{subsection}

% Macros
%
\def\lim{\mathop{\mathrm{lim}}\nolimits}
\def\colim{\mathop{\mathrm{colim}}\nolimits}
\def\Spec{\mathop{\mathrm{Spec}}}
\def\Hom{\mathop{\mathrm{Hom}}\nolimits}
\def\Ext{\mathop{\mathrm{Ext}}\nolimits}
\def\SheafHom{\mathop{\mathcal{H}\!\mathit{om}}\nolimits}
\def\SheafExt{\mathop{\mathcal{E}\!\mathit{xt}}\nolimits}
\def\Sch{\mathit{Sch}}
\def\Mor{\mathop{\mathrm{Mor}}\nolimits}
\def\Ob{\mathop{\mathrm{Ob}}\nolimits}
\def\Sh{\mathop{\mathit{Sh}}\nolimits}
\def\NL{\mathop{N\!L}\nolimits}
\def\CH{\mathop{\mathrm{CH}}\nolimits}
\def\proetale{{pro\text{-}\acute{e}tale}}
\def\etale{{\acute{e}tale}}
\def\QCoh{\mathit{QCoh}}
\def\Ker{\mathop{\mathrm{Ker}}}
\def\Im{\mathop{\mathrm{Im}}}
\def\Coker{\mathop{\mathrm{Coker}}}
\def\Coim{\mathop{\mathrm{Coim}}}

% Boxtimes
%
\DeclareMathSymbol{\boxtimes}{\mathbin}{AMSa}{"02}

%
% Macros for moduli stacks/spaces
%
\def\QCohstack{\mathcal{QC}\!\mathit{oh}}
\def\Cohstack{\mathcal{C}\!\mathit{oh}}
\def\Spacesstack{\mathcal{S}\!\mathit{paces}}
\def\Quotfunctor{\mathrm{Quot}}
\def\Hilbfunctor{\mathrm{Hilb}}
\def\Curvesstack{\mathcal{C}\!\mathit{urves}}
\def\Polarizedstack{\mathcal{P}\!\mathit{olarized}}
\def\Complexesstack{\mathcal{C}\!\mathit{omplexes}}
% \Pic is the operator that assigns to X its picard group, usage \Pic(X)
% \Picardstack_{X/B} denotes the Picard stack of X over B
% \Picardfunctor_{X/B} denotes the Picard functor of X over B
\def\Pic{\mathop{\mathrm{Pic}}\nolimits}
\def\Picardstack{\mathcal{P}\!\mathit{ic}}
\def\Picardfunctor{\mathrm{Pic}}
\def\Deformationcategory{\mathcal{D}\!\mathit{ef}}
\newenvironment{reference}{}{}
\newenvironment{slogan}{}{}
\newenvironment{history}{}{}
\usepackage{cleveref}
\crefname{originallemma}{original lemma}{original lemmas}
\Crefname{originallemma}{Original lemma}{Original lemmas}
\makeatletter
\newlabel{algebra-lemma-integral-no-inclusion}{{00GT}{source}{}{sourcepointer-1}{}}
\newlabel{algebra-lemma-integral-no-inclusion@cref}{{[originallemma][0][]00GT}{[1][0][]source}{}{}{}}
\makeatother
\crefname{originallemma}{original lemma}{original lemmas}
\Crefname{originallemma}{Original lemma}{Original lemmas}
\makeatletter
\newlabel{lemma-stein-universally-closed}{{03GY}{source}{}{sourcepointer-2}{}}
\newlabel{lemma-stein-universally-closed@cref}{{[originallemma][0][]03GY}{[1][0][]source}{}{}{}}
\makeatother
\crefname{originallemma}{original lemma}{original lemmas}
\Crefname{originallemma}{Original lemma}{Original lemmas}
\makeatletter
\newlabel{morphisms-lemma-point-finite-type}{{01TA}{source}{}{sourcepointer-3}{}}
\newlabel{morphisms-lemma-point-finite-type@cref}{{[originallemma][0][]01TA}{[1][0][]source}{}{}{}}
\makeatother
\crefname{originallemma}{original lemma}{original lemmas}
\Crefname{originallemma}{Original lemma}{Original lemmas}
\makeatletter
\newlabel{spaces-morphisms-lemma-enough-finite-type-points}{{06EK}{source}{}{sourcepointer-4}{}}
\newlabel{spaces-morphisms-lemma-enough-finite-type-points@cref}{{[originallemma][0][]06EK}{[1][0][]source}{}{}{}}
\makeatother
\crefname{originallemma}{original lemma}{original lemmas}
\Crefname{originallemma}{Original lemma}{Original lemmas}
\makeatletter
\newlabel{spaces-morphisms-lemma-finite-presentation-noetherian}{{04ZL}{source}{}{sourcepointer-5}{}}
\newlabel{spaces-morphisms-lemma-finite-presentation-noetherian@cref}{{[originallemma][0][]04ZL}{[1][0][]source}{}{}{}}
\makeatother
\title{In the Stein factorization of a universally closed finite-type map of spaces to an affine base, the fibres of pi are finite with finite residue extensions}
\author{The Stacks Project Authors}
\date{stacks-project:spaces-more-morphisms.tex:local-snapshot-20261011; GNU FDL 1.2 or later}
\begin{document}
\maketitle
\noindent\textit{Editorial note.} Statement span is lines 8060-8066 of sources/stacks-spaces-more-morphisms/pinned-original-source.tex; proof span is lines 8068-8088, the complete author proof (19 body lines).\par\medskip
\noindent\textit{Editorial note.} Navigation stubs: 5 cited-but-not-included labels are supplied as non-load-bearing pointers with their pinned tags where the snapshot tag index has them.\par\medskip
\noindent\textit{Editorial note.} extra\_packages is empty on purpose: the retained preamble already loads xy with options and a second declaration clashes in the assembler wrapper.\par\medskip
\noindent\textit{Editorial note.} Text is retained verbatim, including the author's own typos and placeholders.\par\medskip
\noindent\textit{Editorial note.} The author explicitly writes 'A topological argument (omitted)' for the finiteness of the point set; everything else is written out with exact citations.\par\medskip
\noindent\textit{Editorial note.} Counted B side; A (schemes) twin uncounted per ruling.\par\medskip
\noindent\textit{Allowed standard prerequisites (editorial).}\begin{itemize}
\item Explicit prerequisite boundary. In-source lemma-stein-universally-closed; Algebra, lemma integral-no-inclusion; Morphisms of Spaces, lemmas finite-presentation-noetherian and enough-finite-type-points; More on Morphisms, lemma point-finite-type.
\end{itemize}\medskip
\ifdefined\hypertarget\hypertarget{sourcepointer-1}{}\fi
\noindent\textit{Source reference.} Original lemma 00GT: sources/stacks-spaces-more-morphisms/pinned-original-source.tex, source lines 8088--8088. Cited by the retained statement or proof but not part of the collected spans; supplied only as a navigation pointer. The argument does not depend on its content.\par\medskip
\ifdefined\hypertarget\hypertarget{sourcepointer-2}{}\fi
\noindent\textit{Source reference.} Original lemma 03GY: sources/stacks-spaces-more-morphisms/pinned-original-source.tex, source lines 7994--8014. Cited by the retained statement or proof but not part of the collected spans; supplied only as a navigation pointer. The argument does not depend on its content.\par\medskip
\ifdefined\hypertarget\hypertarget{sourcepointer-3}{}\fi
\noindent\textit{Source reference.} Original lemma 01TA: sources/stacks-spaces-more-morphisms/pinned-original-source.tex, source lines 8088--8088. Cited by the retained statement or proof but not part of the collected spans; supplied only as a navigation pointer. The argument does not depend on its content.\par\medskip
\ifdefined\hypertarget\hypertarget{sourcepointer-4}{}\fi
\noindent\textit{Source reference.} Original lemma 06EK: sources/stacks-spaces-more-morphisms/pinned-original-source.tex, source lines 8088--8088. Cited by the retained statement or proof but not part of the collected spans; supplied only as a navigation pointer. The argument does not depend on its content.\par\medskip
\ifdefined\hypertarget\hypertarget{sourcepointer-5}{}\fi
\noindent\textit{Source reference.} Original lemma 04ZL: sources/stacks-spaces-more-morphisms/pinned-original-source.tex, source lines 8088--8088. Cited by the retained statement or proof but not part of the collected spans; supplied only as a navigation pointer. The argument does not depend on its content.\par\medskip
\begin{lemma}
\label{lemma-stein-universally-closed-residue-fields}
In Lemma \ref{lemma-stein-universally-closed} assume in addition that
$f$ is locally of finite type and $Y$ affine. Then for $y \in Y$ the fibre
$\pi^{-1}(\{y\}) = \{y_1, \ldots, y_n\}$ is finite and the field extensions
$\kappa(y_i)/\kappa(y)$ are finite.
\end{lemma}

\begin{proof}
Recall that there are no specializations among the points of $\pi^{-1}(\{y\})$,
see Algebra, Lemma \ref{algebra-lemma-integral-no-inclusion}.
As $f'$ is surjective, we find that $|X_y| \to \pi^{-1}(\{y\})$ is surjective.
Observe that $X_y$ is a quasi-separated algebraic space of finite type
over a field (quasi-compactness was shown in the proof of the
referenced lemma). Thus $|X_y|$ is a Noetherian topological space
(Morphisms of Spaces, Lemma
\ref{spaces-morphisms-lemma-finite-presentation-noetherian}).
A topological argument (omitted) now shows that $\pi^{-1}(\{y\})$ is finite.
For each $i$ we can pick a finite type point $x_i \in |X_y|$
mapping to $y_i$ (Morphisms of Spaces, Lemma
\ref{spaces-morphisms-lemma-enough-finite-type-points}).
We conclude that $\kappa(y_i)/\kappa(y)$ is finite:
$x_i$ can be represented by a morphism $\Spec(k_i) \to X_y$
of finite type (by our definition of finite type points)
and hence $\Spec(k_i) \to y = \Spec(\kappa(y))$ is of finite type
(as a composition of finite type morphisms),
hence $k_i/\kappa(y)$ is finite (Morphisms, Lemma
\ref{morphisms-lemma-point-finite-type}).
\end{proof}

\end{document}
